---
format: yohaku/0.2
target: text
macro: '[[macros/raw]]'
---

## 正文
````text
//! X11 canvas reuses the IME probe's native drawing path and existing font fallback.
use std::ffi::{c_int, c_uint, c_ulong, c_void};
use winit::{
    raw_window_handle::{HasDisplayHandle, HasWindowHandle, RawDisplayHandle, RawWindowHandle},
    window::Window,
};
#[link(name = "X11")]
unsafe extern "C" {
    fn XDefaultGC(d: *mut c_void, s: c_int) -> *mut c_void;
    fn XDefaultVisual(d: *mut c_void, s: c_int) -> *mut c_void;
    fn XDefaultDepth(d: *mut c_void, s: c_int) -> c_int;
    fn XCreateImage(
        d: *mut c_void,
        v: *mut c_void,
        depth: c_uint,
        format: c_int,
        offset: c_int,
        data: *mut i8,
        w: c_uint,
        h: c_uint,
        pad: c_int,
        stride: c_int,
    ) -> *mut c_void;
    fn XPutImage(
        d: *mut c_void,
        w: c_ulong,
        gc: *mut c_void,
        image: *mut c_void,
        sx: c_int,
        sy: c_int,
        dx: c_int,
        dy: c_int,
        width: c_uint,
        height: c_uint,
    ) -> c_int;
    fn XDestroyImage(image: *mut c_void) -> c_int;
    fn XFlush(d: *mut c_void) -> c_int;
    fn malloc(size: usize) -> *mut c_void;
    fn free(p: *mut c_void);
}
pub struct Canvas {
    pub width: usize,
    pub height: usize,
    pub pixels: Vec<u32>,
}
impl Canvas {
    pub fn new(width: usize, height: usize) -> Self {
        Self {
            width,
            height,
            pixels: vec![0x101820; width * height],
        }
    }
    pub fn rect(&mut self, x: i32, y: i32, w: i32, h: i32, color: u32) {
        for row in y.max(0)..(y + h).min(self.height as i32) {
            let start = row as usize * self.width + x.max(0).min(self.width as i32) as usize;
            let end = row as usize * self.width + (x + w).max(0).min(self.width as i32) as usize;
            if start <= end {
                self.pixels[start..end].fill(color)
            }
        }
    }
    pub fn text(&mut self, text: &str, x: i32, y: i32, color: u32, max: i32) {
        let mut pen = x;
        for ch in text.chars() {
            if pen >= max {
                break;
            }
            let (metrics, bitmap) = super::text_font::for_character(ch).rasterize(ch, 18.);
            for row in 0..metrics.height {
                for col in 0..metrics.width {
                    let alpha = bitmap[row * metrics.width + col] as u32;
                    if alpha == 0 {
                        continue;
                    }
                    let px = pen + metrics.xmin + col as i32;
                    let py = y + 19 - metrics.ymin - metrics.height as i32 + row as i32;
                    if px >= 0
                        && py >= 0
                        && (px as usize) < self.width
                        && (py as usize) < self.height
                    {
                        let at = py as usize * self.width + px as usize;
                        let old = self.pixels[at];
                        let mix = |shift: u32| {
                            ((((color >> shift) & 255u32) * alpha
                                + ((old >> shift) & 255u32) * (255 - alpha))
                                / 255)
                                << shift
                        };
                        self.pixels[at] = mix(16) | mix(8) | mix(0);
                    }
                }
            }
            pen += 10 * super::terminal_screen::width(ch) as i32;
        }
    }
    pub fn present(&self, window: &Window) -> Result<(), String> {
        let (Ok(display), Ok(handle)) = (window.display_handle(), window.window_handle()) else {
            return Err("Window handles unavailable".into());
        };
        let (RawDisplayHandle::Xlib(display), RawWindowHandle::Xlib(handle)) =
            (display.as_raw(), handle.as_raw())
        else {
            return Err(
                "This minimal terminal window requires the existing X11 host display".into(),
            );
        };
        let Some(d) = display.display else {
            return Err("X11 display unavailable".into());
        };
        window.pre_present_notify();
        unsafe {
            let size = self.pixels.len() * 4;
            let data = malloc(size);
            if data.is_null() {
                return Err("Canvas allocation failed".into());
            }
            std::ptr::copy_nonoverlapping(
                self.pixels.as_ptr().cast::<u8>(),
                data.cast::<u8>(),
                size,
            );
            let image = XCreateImage(
                d.as_ptr(),
                XDefaultVisual(d.as_ptr(), display.screen),
                XDefaultDepth(d.as_ptr(), display.screen) as u32,
                2,
                0,
                data.cast(),
                self.width as u32,
                self.height as u32,
                32,
                (self.width * 4) as i32,
            );
            if image.is_null() {
                free(data);
                return Err("X11 image creation failed".into());
            }
            XPutImage(
                d.as_ptr(),
                handle.window,
                XDefaultGC(d.as_ptr(), display.screen),
                image,
                0,
                0,
                0,
                0,
                self.width as u32,
                self.height as u32,
            );
            XDestroyImage(image);
            XFlush(d.as_ptr());
        }
        Ok(())
    }
}
````

## 输出与关联
```yaml
output: "examples/yohaku_workspace/terminal_canvas.rs"
```

## AI 笔记
```yaml
{"source_version": 1, "writer": "Yohaku release packaging", "summary": "Portable release source. Scope and packaging adjustments are documented in PROVENANCE.md."}
```
